%% kapa: exam=1 round=35 year=2024 period=1 form=A subject=civil_law subject_ko=민법 seq=28 number=28
%% source: Q-Net 감정평가사 기출문제 (cst00309) artlSeq=5213229 file=raw/1008/5213229/1교시 문제지 원본.pdf
%% conversion: kordoc 4.24.51 -> tools/gen1.py (auto); review=proofread-sonnet; flags=none
\begin{question}{28}
부합에 관한 설명으로 옳지 않은 것은? (다툼이 있으면 판례에 따름)
\begin{choices}
  \item 부동산에 부합되어 동산의 소유권이 소멸한 때에는 그 동산을 목적으로 한 다른 권리도 소멸한다.
  \item 부합한 동산 간의 주종을 구별할 수 없는 때에는 특약이 없는 한 동산의 소유자는 부합당시 가액의 비율로 합성물을 공유한다.
  \item X토지 소유자의 승낙없이 토지임차인의 승낙만 받아 제3자가 X에 수목을 심은 경우, 그 수목은 X에 부합하지 않으므로 제3자가 식재한 수목임을 알지 못하는 X의 양수인은 그 수목을 벌채할 수 없다.
  \item 타인의 권원에 기하여 부동산에 부합된 물건이 부동산의 구성부분이 된 경우, 부동산의 소유자는 방해배제청구권에 기하여 부합물의 철거를 청구할 수 없다.
  \item 건물의 증축부분이 축조 당시 독립한 권리의 객체성을 상실하여 본건물에 부합된 후 구조의 변경 등으로 독립한 권리의 객체성을 취득하게 된 때에는 본건물과 독립하여 거래의 대상이 될 수 있다.
\end{choices}
\end{question}
\answer{3}
